\begin{frame}[shrink]
	\frametitle{Codificação Poisson: intensidade vira chance}
	\small
	\par Uma rede precisa receber números comuns, mas seus neurônios trocam pulsos. Na codificação Poisson, o valor de entrada não determina um pulso certo: determina a \textbf{probabilidade} de disparar em cada passo.
	\begin{center}
		\begin{tabular}{lcl}
			\toprule
			\textbf{Entrada} & \textbf{Probabilidade por passo} & \textbf{Leitura após 10 passos} \\
			\midrule
			$x=0{,}6$, $r_{\max}=0{,}9$ & $p=x\,r_{\max}=0{,}54$ & $E[n]=10p=5{,}4$ pulsos \\
			\bottomrule
		\end{tabular}
	\end{center}
	\begin{itemize}
		\item Duas repetições com a mesma entrada podem gerar contagens diferentes; ambas tendem a representar a mesma intensidade.
		\item \textbf{Por que importa?} A contagem aproxima o valor, mas traz variação aleatória. Mais passos tornam a estimativa de taxa mais estável, ao custo de mais tempo de simulação.
		\item É uma escolha de representação, não ruído adicionado depois do treino. A codificação por taxa é comum, mas não é a única forma de representar informação em pulsos \cite{diehl2015fast, eshraghian2023training}.
	\end{itemize}
	\begin{block}{O que observar na demonstração}
		O mesmo sinal alimenta cada passo. A intensidade define a chance de disparo; não garante que um pulso apareça em um instante específico.
	\end{block}
\end{frame}
